% !TEX root = ../thesis.tex

配备机械臂的腿式机器人能够执行单纯移动平台或固定底座操作系统均无法独立完成的复杂移动操作任务。
本项目将上海交通大学开源的 \textit{RoboDuet} 框架——最初为 Unitree Go1+Z1 平台设计——迁移至更重、
功能更强的 \textit{Unitree B2+Z1} 系统。该系统由重约 60\,kg 的四足机器人本体与一个六自由度机械臂组成。

本项目的核心技术贡献在于完整迁移了去中心化多智能体强化学习（MARL）训练流水线，涵盖机器人 URDF 集成、
自由度（DOF）映射、PD 增益标定，以及基于 NVIDIA IsaacGym 的仿真演示框架搭建。
系统采用近端策略优化（PPO）算法训练两个解耦策略：
一个\textit{运动策略}，用于将四足机器人稳定在标准站立高度（$h\approx0.45$\,m）；
另一个\textit{机械臂策略}，采用基于历史观测的 Actor-Critic 架构，通过球形坐标末端执行器指令
$(l, p, \psi)$（即伸展距离、俯仰角与偏航角）实现跟踪，并配备可学习的自适应模块以支持仿真到真实的迁移。

在仿真环境中共完成三项演示：
（1）固定底座机械臂操作，验证末端执行器向前伸展方向的正确性；
（2）四足机器人站立状态下同时进行机械臂伸展，验证同步控制的稳定性；
（3）固定底座抓取与提升序列：机械臂到达一个边长 6\,cm 的小方块处，夹爪闭合后将方块提升 30\,cm，
通过运动学方块瞬移策略完整展示"到达—抓取—提升"流程。

演示二中机器人底座高度始终保持在标称值 $\pm$3\,cm 以内；
演示一与演示三中底座以运动学固定方式模拟，高度恒为 0.340\,m，
验证了去中心化双策略架构在 B2Z1 平台上的稳定性，无需任何再训练。

上述结果确认了原始 RoboDuet 论文中关于跨形态零样本迁移能力的声明，
并表明 B2Z1 平台是未来真实世界移动操作实验的可行测试平台。
